\documentclass[../ICIT-Paper-2024.tex]{subfiles}
\begin{document}
\paragraph{}
Stochastic gradient descent (SGD) is an iterative optimization algorithm used in machine learning to find the values of parameters (such as weights) that minimize a loss function. It is a variation of the standard gradient descent algorithm.

Stochastic Gradient Descent (SGD) is a popular algorithm that can achieve state-of-the-art performance on a variety of machine learning tasks~\cite{recht2011hogwild}. In standard gradient descent, the gradient is computed using the entire training dataset. In SGD, the gradient is computed using only a single randomly selected training example (or a small batch of examples) at each iteration. This reduces the computational cost per iteration. The pseudocode~\ref{al1} describes in detail our algorithm.
\begin{algorithm}[H]
\caption{Stochastic Gradient Descent Algorithm Implementation}\label{al1}
\label{alg:example}
\begin{algorithmic}[1]
\Require Learning rate $\alpha$, decay rate $decay\_rate$, num\_iterations
\State Initialize parameter $\theta$
\State Time step $t = 0$
\For{$i$ in range(num\_iterations)}
    \State Sample a mini-batch of $m$ training examples $x_i$ with corresponding labels $\mathrm{first\_Y}_i$, $\mathrm{second\_Y}_i$
    \State Compute gradients: $g \leftarrow \frac{1}{m} \nabla_\theta \sum_{i=1}^m J(\theta; x_i, \mathrm{first\_Y}_i, \mathrm{second\_Y}_i)$
    \State \textbf{Update parameters:} $\theta \leftarrow \theta - \alpha g$
    \State $\alpha \leftarrow \alpha * decay\_rate$
    \State $t \leftarrow t + 1$
\EndFor
\end{algorithmic}
\end{algorithm}


In the SGD application section, we adjusted the loss function to fit the problem of assigning 2 labels. With the input parameters, $\mathrm{first\_Y}$ representing the first label, $\mathrm{second\_Y}$ representing the second label, and $\mathrm{Y\_pred}$ representing the predicted value, the new loss function is defined as the formula~\ref{eq2}:
\begin{equation}\label{eq2}
E(i) = \left\{
\begin{array}{ll}
y\_i^p - y\_i^f & \text{if } y\_i^f = y\_i^s \\
y\_i^p - min(y\_i^f, y\_i^s) & \text{if } y\_i^p < min(y\_i^f, y\_i^s) \\
y\_i^p - max(y\_i^f, y\_i^s) & \text{if } y\_i^p > max(y\_i^f, y\_i^s)
\end{array}
\right.
\end{equation}

\end{document}